\chapter{Diseño y arquitectura}

\section{Visión general}
La arquitectura del framework se organiza en dos artefactos desplegables
--- el backend \emph{uniback} y el frontend \emph{ngt-gui} --- unidos por un
\textbf{contrato} explícito que actúa como única superficie de acoplamiento.
El principio rector del diseño es que el \emph{modelo de datos} (los modelos
ORM de SQLAlchemy) es la única fuente de verdad: de él se derivan, de forma
automática, el esquema, el contrato de la API y los formularios de la
interfaz. El Algoritmo~\ref{lst:arq} representa el flujo de datos completo.

\begin{lstlisting}[style=docker, caption={Flujo de datos extremo a extremo del framework}, label={lst:arq}]
 Modelo SQLAlchemy (fuente de verdad)
        |  make_simple_rest_crud / make_crudie_rest_crud
        v
 sqlalchemy_to_pydantic  --> Pydantic + hints x-* por columna
        |
        v
 get_enriched_json_schema --> JSON Schema enriquecido
        |
        +--> GET /<entity>/schema.json  (ETag)
        +--> GET /api/sys/schemas        (bundle versionado)
        +--> GET /api/sys/entities       (descubrimiento)
        |
        v   [ ResponseEnvelope: content / count / issues ]
 ----------------- contrato HTTP -----------------
        |
        v
 ngt-gui: EntityClient / SchemaService / DiscoveryService
        |
        v
 schemaToFormly(schema)  --> FormlyFieldConfig[]
        |
        v
 <ngt-dynamic-form>  --> formulario renderizado
\end{lstlisting}

\section{Arquitectura del backend (\emph{uniback})}
El backend sigue una arquitectura en capas dentro de un paquete Python
instalable, construido sobre FastAPI~\cite{fastapi} y
SQLAlchemy~2.x~\cite{sqlalchemy}. Las capas principales son:

\begin{description}
  \item[Persistencia.] Modelos ORM SQLAlchemy~2.x organizados por dominio
  (\texttt{core}, \texttt{sysadmin}, \texttt{annotations},
  \texttt{hierarchies}, \texttt{species}, \texttt{files},
  \texttt{geographics}, \texttt{screens}). Un \texttt{BaseMixin} aporta
  funcionalidad común (UUID, marcas de tiempo, atributos). El prefijo de
  tablas es configurable (\texttt{ub\_} por defecto) y \gls{continuum}
  proporciona el histórico de cambios.
  \item[Servicios.] Lógica de dominio (autenticación, colecciones,
  jerarquías, importación, anotaciones) desacoplada de la capa HTTP.
  \item[API.] FastAPI: \emph{routers}, esquemas Pydantic, manejadores de
  error globales, fábricas CRUD y registro de entidades.
  \item[Inicializador.] Un único punto de entrada \texttt{initialize(config)}
  que normaliza la configuración (Pydantic Settings), prepara la base ORM, el
  gestor de sesiones y la aplicación FastAPI, e inicializa de forma opcional
  Redis, Celery y Firebase.
\end{description}

El despliegue se realiza con \gls{docker} mediante \texttt{docker compose},
con servicios para PostgreSQL, Redis, el \emph{worker} Celery y el servidor
web, lo que garantiza reproducibilidad multiplataforma (RNF-02). Existe
además un \texttt{api\_gateway} con configuración Nginx para escenarios de
microservicios y un mecanismo de descubrimiento de servicios externos a
través de Traefik.

\section{Sistema de plugins y extensibilidad}\label{sec:plugins}
El framework no debe acoplar ninguna lógica a un dominio concreto (RNF-01); el
dominio se introduce \emph{desde fuera} del núcleo mediante un sistema de
\gls{plugin}s. Un plugin es un módulo Python autónomo que se engancha al ciclo
de vida del inicializador para aportar modelos, \emph{endpoints}, datos semilla
y extensiones de comportamiento, sin que el núcleo conozca su existencia. Es la
pieza que materializa la idea, heredada de la transición desde
\gls{nextgendem}, de que el caso de uso del \gls{jbvc} sea una capa apilada
sobre \emph{uniback} y no una modificación de él.

\subsection{El contrato del plugin}
Todo plugin hereda de la clase base \texttt{UnibackPlugin} y sobreescribe solo
los \emph{hooks} que necesita; el resto devuelven listas vacías o no hacen nada,
lo que preserva la retrocompatibilidad. Los \emph{hooks} se dividen en eventos
de ciclo de vida y \emph{getters} de extensiones inyectables, resumidos en la
tabla~\ref{tab:plugin-hooks}.

\begin{table}[!htbp]
\caption{\emph{Hooks} de ciclo de vida y extensiones inyectables de \texttt{UnibackPlugin}.}
\label{tab:plugin-hooks}
\centering
\begin{tabular}{|l|l|p{6.3cm}|}
\hline
\textbf{Hook} & \textbf{Tipo} & \textbf{Propósito} \\ \hline
\texttt{on\_init(settings)} & ciclo de vida & Configuración temprana y
servicios de terceros. \\ \hline
\texttt{get\_model\_modules()} & ciclo de vida & Módulos con modelos ORM a
registrar antes de \texttt{configure\_mappers()}. \\ \hline
\texttt{get\_routers()} & ciclo de vida & \emph{Routers} FastAPI a montar en la
aplicación. \\ \hline
\texttt{on\_seed(db)} & ciclo de vida & Inyectar filas por defecto durante el
\emph{auto-seed}. \\ \hline
\texttt{on\_app\_ready(app)} & ciclo de vida & Middlewares o montajes tras crear
la app. \\ \hline
\texttt{get\_source\_adapters()} & extensión & Adaptadores que abren
\emph{streams} desde orígenes externos. \\ \hline
\texttt{get\_importers()}, \texttt{get\_exporters()} & extensión & Parsers y
serializadores para el insertador genérico. \\ \hline
\texttt{get\_field\_widgets()} & extensión & Pistas de UI que enriquecen el
JSON Schema generado. \\ \hline
\texttt{get\_fk\_resolvers()} & extensión & Resolvers de claves foráneas
(\texttt{\{by, value\}}) en la importación. \\ \hline
\end{tabular}
\end{table}

\subsection{Descubrimiento e integración en el inicializador}
El \texttt{PluginManager} (un \emph{singleton} global) descubre los plugins
escaneando un directorio configurable (variable \texttt{UNIBACK\_PLUGINS\_PATH},
por defecto \texttt{src/plugins}): por cada paquete importa su submódulo
\texttt{plugin} y registra las subclases de \texttt{UnibackPlugin} que encuentra.
A partir de ahí, \texttt{initialize(config)} emite los eventos del ciclo de vida
en un orden preciso, intercalándolos con la construcción de la base ORM y de la
aplicación FastAPI. El Algoritmo~\ref{lst:plugin-lifecycle} resume esa
secuencia.

\begin{lstlisting}[style=docker, caption={Integración de los plugins en el inicializador}, label={lst:plugin-lifecycle}]
 initialize(config)
   |
   v
 discover_plugins(UNIBACK_PLUGINS_PATH | src/plugins)
   |   importa <plugin>.plugin y detecta subclases de UnibackPlugin
   v
 emit_on_init(settings)        # hook on_init de cada plugin
 register_core_defaults()      # importers/resolvers del nucleo
 populate_registries()         # vuelca los get_* en los registros globales
   |
   v
 create_orm_base() + import modelos del nucleo y de los plugins  # get_model_modules
   |
   v
 configure_mappers()           # polimorfismo + versionado (Continuum)
   |
   v
 create_tables()  (si procede)
   |
   v
 auto_seed: emit_on_seed(db)   # hook on_seed dentro del session_scope
   |
   v
 create_app(...)               # monta los get_routers() de los plugins
   |
   v
 emit_on_app_ready(app)        # hook on_app_ready
\end{lstlisting}

El orden es deliberado: los modelos del plugin
(\texttt{get\_model\_modules}) se importan \emph{antes} de
\texttt{configure\_mappers()}, de modo que sus entidades queden registradas en
la misma base ORM que las del núcleo y participen del polimorfismo y del
versionado; \texttt{on\_seed} corre dentro del \emph{scope} de sesión del
inicializador, junto a la siembra del núcleo; y \texttt{on\_app\_ready} se
invoca con la aplicación ya construida, momento en el que el plugin puede montar
middlewares o ficheros estáticos.

\subsection{Extensiones inyectables y registros}
Más allá de los \emph{hooks} de ciclo de vida, un plugin puede aportar
implementaciones de cinco contratos (\texttt{SourceAdapter}, \texttt{Importer},
\texttt{Exporter}, \texttt{FieldWidget} y \texttt{FKResolver}) que el
\texttt{PluginManager} vuelca en sendos registros globales \emph{thread-safe}.
La política de selección es ``el último registrado para un mismo \texttt{name}
gana'', de modo que un plugin puede sobreescribir un comportamiento por defecto
del núcleo simplemente registrando otra implementación con el mismo nombre. El
núcleo inscribe primero sus \emph{defaults} (\texttt{JsonImporter},
\texttt{NdjsonImporter} y dos resolvers de claves foráneas), de manera que la
importación genérica funciona \emph{out of the box} sin ningún plugin presente.
Estos contratos son los puntos de extensión que permiten, por ejemplo,
introducir un nuevo formato de importación (CSV, XLSX) o un tipo de \emph{widget}
de UI (un mapa, un editor de código) sin tocar el conversor de esquema ni las
fábricas CRUD.

\subsection{Ejemplo: el plugin de plantas}
El plugin \texttt{PlantsApp} es una aplicación de demostración que añade, como
capa externa, un catálogo de plantas geolocalizadas --- un dominio cercano al
del \gls{jbvc}. Sin modificar una sola línea del núcleo, el plugin
(i)~registra el modelo \texttt{Plant} (subclase de \texttt{FunctionalObject});
(ii)~expone su CRUD REST completo con \texttt{make\_simple\_rest\_crud},
incluyendo \texttt{/schema.json} y \texttt{/bulk}; (iii)~aporta un
\texttt{FieldWidget} que enriquece las columnas de coordenadas con pistas
\texttt{x-*} para el mapa; y (iv)~siembra datos de ejemplo y una entrada en el
menú lateral. El Algoritmo~\ref{lst:plants-plugin} muestra su esqueleto.

\begin{lstlisting}[language=Python, caption={Esqueleto del plugin de plantas (\texttt{plants/plugin.py})}, label={lst:plants-plugin}]
class PlantsPlugin(UnibackPlugin):
    name = "PlantsApp"
    version = "1.0.0"

    # 1) Modelos: registrados antes de configure_mappers()
    def get_model_modules(self):
        return ["plants.models"]

    # 2) API: CRUD REST completo + /schema.json + /bulk
    def get_routers(self):
        from uniback.api.crud_factory import make_simple_rest_crud
        from plants.models import Plant
        return [make_simple_rest_crud(Plant, "plants", tags=["Plants"])]

    # 3) Pista de UI inyectable para las coordenadas del mapa
    def get_field_widgets(self):
        return [MapCoordinatesWidget()]

    # 4) Semilla: ObjectType, datos demo y menu lateral
    def on_seed(self, db):
        ...
\end{lstlisting}

Este ejemplo cierra el argumento de reutilización (RNF-01): exponer un dominio
nuevo y editable desde el frontend se reduce a escribir un plugin que declara su
modelo y lo expone con una fábrica CRUD; el contrato, el esquema enriquecido y
el formulario dinámico se derivan automáticamente, tal como describe el resto
del capítulo.

\section{Arquitectura del frontend (\emph{ngt-gui})}
El frontend es una librería Angular~\cite{angular} publicable
(\texttt{projects/ngt-gui}) acompañada de una aplicación de demostración, que
emplea \gls{formly}~\cite{formly} para construir los formularios a partir del
\gls{jsonschema}~\cite{jsonschema}. La librería se estructura en
\texttt{core} (lógica sin dependencia de UI: conversor de esquema, tipos del
contrato), \texttt{services} (clientes HTTP, caché de esquemas,
descubrimiento), \texttt{components} (catálogo de presentación, entre ellos
\texttt{dynamic-form} y \texttt{formly-components}), \texttt{models},
\texttt{interfaces} y \texttt{tokens} de inyección.

\subsection{Catálogo de modelos del núcleo}
El backend agrupa los modelos por dominio y les asigna un código de tipo de
objeto (\texttt{ObjectType}) que habilita el tratamiento polimórfico. La
tabla~\ref{tab:modelos} resume los grupos principales y su finalidad.

\begin{table}[!htbp]
\caption{Grupos de modelos del núcleo de \emph{uniback}.}
\label{tab:modelos}
\centering
\begin{tabular}{|l|p{8.5cm}|}
\hline
\textbf{Módulo} & \textbf{Contenido} \\ \hline
\texttt{core} & \texttt{ObjectType}, \texttt{FunctionalObject}: base común
con UUID, propiedad, marcas de tiempo y atributos. \\ \hline
\texttt{sysadmin} & Autenticación/autorización: \texttt{Authorizable},
\texttt{Identity}, \texttt{Role}, \texttt{Organization}, \texttt{Group},
\texttt{ACL}. \\ \hline
\texttt{annotations} & Plantillas, campos y textos de anotación sobre objetos
funcionales. \\ \hline
\texttt{hierarchies} & Jerarquías (p. ej. taxonómicas) y colecciones. \\ \hline
\texttt{species} & Modelos de especies (incl. soporte Darwin Core). \\ \hline
\texttt{files} & Sistema de ficheros: \texttt{FileSystemObject},
\texttt{Folder}, \texttt{File}. \\ \hline
\texttt{screens} & Modelos de GUI: \emph{app flavor}, menús, pantallas. \\ \hline
\texttt{geographics} & Datos geográficos asociados a las entidades. \\ \hline
\end{tabular}
\end{table}

\subsection{Catálogo de la librería de componentes}
La librería \texttt{ngt-gui} se organiza en módulos funcionales: \texttt{core}
(conversor de esquema y tipos del contrato, sin dependencia de UI),
\texttt{services} (\texttt{EntityClient}, \texttt{SchemaService},
\texttt{DiscoveryService}, interceptor de error), \texttt{components}
(\texttt{dynamic-form}, \texttt{formly-components} con \textasciitilde20 tipos
Formly personalizados sobre \emph{ng-zorro}, además de \emph{layout},
\emph{data}, \emph{display}, \emph{admin} y \emph{auth}), \texttt{models},
\texttt{interfaces}, \texttt{pipes} y \texttt{tokens} de inyección. Esta
separación garantiza que la lógica de generación (\texttt{core}) sea
testeable de forma aislada (RNF-04).

\section{Persistencia, sesiones y configuración}
La capa de persistencia se apoya en una base ORM declarativa creada por
\texttt{create\_orm\_base()} y en un gestor de sesiones
(\texttt{SessionManager} / \texttt{create\_session\_factory}) que expone la
sesión activa mediante un contexto (\texttt{current\_session\_ctx}). Esto
permite que los modelos --- por ejemplo, para resolver el propietario actual
de un \texttt{FunctionalObject} --- accedan a la sesión sin recibirla por
parámetro, manteniendo el código de dominio limpio.

La configuración se modela con Pydantic Settings (\texttt{Settings},
\texttt{AuthSettings}, \texttt{WorkerSettings}), de modo que un diccionario o
variables de entorno se validan y normalizan en un único punto. El
inicializador acepta indistintamente un diccionario o una instancia
\texttt{Settings}, lo que facilita tanto el uso programático como el despliegue
en contenedores con variables \texttt{UNIBACK\_*}.

\section{Modelo polimórfico de objetos funcionales}
El corazón del modelo de datos es la jerarquía de \texttt{FunctionalObject}.
Cada entidad de dominio hereda de ella y queda asociada a un
\texttt{ObjectType} con un código numérico (p. ej. \emph{dataset}=0,
\emph{collection}=106, \emph{screen}=32). Este diseño habilita el tratamiento
homogéneo de entidades heterogéneas: anotaciones, ACL y colecciones pueden
referirse a ``cualquier objeto funcional'' sin conocer su clase concreta. El
Algoritmo~\ref{lst:erd} resume la relación de forma textual.

\begin{lstlisting}[style=docker, caption={Esquema conceptual (textual) del núcleo}, label={lst:erd}]
 ObjectType (id, uuid, name)
     ^ 1
     | N
 FunctionalObject (id, uuid, owner, created/updated, attrs)
     ^                         ^
     | hereda                  | referido por
 [Collection, Screen, ...]  Annotation* / ACL*

 Authorizable (name, uuid)        <- tabla padre
     ^                ^
   Identity         Role           <- herencia polimórfica
\end{lstlisting}

\section{Modelo de autorización}
La autorización se basa en \emph{Authorizables} (identidades y roles, con
herencia polimórfica), organizaciones y grupos, y un modelo de permisos
formado por \texttt{ACL}, \texttt{ACLExpression} y \texttt{ACLDetail} junto con
\texttt{PermissionType} y \texttt{ObjectTypePermissionType}. Este diseño
permite expresar permisos por tipo de objeto y por instancia, soportando el
principio de mínimo privilegio. Existen roles e identidades de sistema con
UUID fijos (p. ej. \emph{sys-admin}, \emph{guest}) para arranque y pruebas. La
autenticación admite Firebase (verificación de \emph{token}) y un modo local
de desarrollo, claramente separado para no exponerse en producción (RNF-07).

\section{El contrato como frontera de desacoplamiento}
La decisión de diseño más relevante es interponer un \textbf{envoltorio
universal} (\emph{ResponseEnvelope}) entre ambos lados. Todo el tráfico, de
éxito o de error, viaja con la misma forma --- \texttt{content}, \texttt{count}
e \texttt{issues} --- y un catálogo de códigos de error estándar. Gracias a
ello, el frontend nunca necesita conocer la forma concreta de una entidad para
manejar respuestas y errores: solo el conversor de esquema, en un único punto,
traduce datos a interfaz. Esta frontera es lo que hace posible cumplir los
requisitos de reutilización (RNF-01) y mantenibilidad (RNF-03).

\section{Escenario de integración extremo a extremo}
El contrato completo se ejercita componiendo las piezas reutilizables de
\texttt{ngt-gui}; en particular, el componente \texttt{<ngt-dynamic-form>}
resuelve internamente el descubrimiento, la obtención y cacheo del esquema, la
conversión a Formly y el envío mediante \texttt{EntityClient}. El
Algoritmo~\ref{lst:seq} describe la secuencia de interacción resultante, que se
verifica con las pruebas automatizadas del conversor y del bundle.

\begin{lstlisting}[style=docker, caption={Secuencia de integración extremo a extremo}, label={lst:seq}]
 Vista            DiscoveryService   SchemaService    EntityClient   Backend
   |  loadEntities()   |                  |                |            |
   |------------------>|---- GET /api/sys/entities ------------------->|
   |   EntityDescriptor[] <-------------------------------------------|
   |  selecciona "roles"|                  |                |          |
   |---- get('/roles') ------------------->| GET /roles/schema.json -->|
   |   JSON Schema (cache) <---------------|<-- ETag / 304 ------------|
   |  schemaToFormly(schema) -> fields     |                |          |
   |  usuario rellena y envia              |                |          |
   |---- create(model) ------------------------------------>| POST /roles/ ->|
   |   ResponseEnvelope (content/issues) <------------------|<---------|
\end{lstlisting}

Este escenario demuestra que las cuatro piezas --- descubrimiento, esquema
cacheado con ETag, conversión a Formly y cliente tipado --- encajan sin
acoplamiento al dominio: cambiar de \texttt{roles} a cualquier otra entidad no
requiere tocar una sola línea de la vista que monta el formulario.
